\subsection{Related Work}\label{sec:relwork}

GANs have been applied to financial time series beyond the ForGAN and Fin-GAN framework described in Sec.~\ref{sec:gans}. Three directions are relevant for this thesis. The first direction focuses on alternative GAN architectures and cost functions for generating realistic financial time series. The second direction investigates GAN-based simulators of option markets. The third direction examines applications of GANs in energy markets.  

In the first direction, \textcite{yoonTimeseriesGenerativeAdversarial} introduce TimeGAN, which trains the generator with both a cross-entropy loss and a supervised next-step prediction loss computed in an embedding space. Quant GANs keep the standard cross-entropy loss and replace the recurrent architecture used in earlier time-series GANs with temporal convolutional networks to capture volatility clustering \parencite{wieseQuantGANsDeep2020}. \textcite{niConditionalSigWassersteinGANs2020} propose the conditional Sig-Wasserstein GAN, which replaces the trained discriminator with an explicitly computed distance between summary statistics of whole paths. \textcite{zhangVisualRealismReliable2026} introduce the Stylized Facts Alignment GAN, which adds stylized facts as differentiable constraints to the cost function of the generator. 

In the second direction, \textcite{wieseDeepHedgingLearning2019} build a GAN-based simulator of an equity option market and show that adversarial training reproduces the historical skewness, kurtosis, and autocorrelation of implied volatilities more accurately than a VAR model or quasi-maximum-likelihood estimation. \textcite{vuleticVolGANGenerativeModel2024} introduce VolGAN, a conditional GAN that simulates the return of the underlying asset and its implied volatility surface. VolGAN adds a smoothness penalty to the cost function of the generator so that the simulated surfaces avoid unrealistic jumps between neighboring moneyness levels and expiries. 

In the third direction, \textcite{bohorquezcorreaTimevaryingSystemicRisk2026} apply TimeGAN to 25 energy markets, including natural gas, to generate synthetic price data and use it to estimate crisis-risk indicators. \textcite{walterProbabilisticSimulationElectricity2024} propose a conditional GAN architecture combining convolutional and recurrent layers to generate day-ahead electricity price scenarios for the German market. \textcite{kornSimulatingIntradayElectricity2025} apply the unmodified ForGAN architecture with a GRU cell to German electricity consumption data and show that it reproduces the recurring daily and weekly patterns in the series.

Across these three directions, GANs have been used to generate realistic financial time series, to simulate option markets, and to model energy market data. However, no study uses GAN-generated paths of natural gas futures to price options. The proposed GAN architectures and cost functions are developed and tested on equity or general financial data. The option-market simulators are restricted to equity options, and the energy market applications generate synthetic consumption or price data for forecasting and risk monitoring. This thesis addresses this gap by applying the ForGAN and Fin-GAN frameworks to TTF natural gas futures data and evaluating the generated paths through option prices.